\documentclass[prl,twocolumn,superscriptaddress,amsmath,amssymb]{revtex4-2}

\usepackage{graphicx}
\usepackage{bm}
\usepackage{hyperref}

\begin{document}

\title{Supplemental Material:\\Ergotropy Rate Equation and Optimal Power Extraction\\from Quantum Non-Equilibrium Steady States}

\author{Zhongchang Huang}
\email{valhuang@kaiwucl.com}
\affiliation{Independent Researcher}

\date{\today}

\maketitle

\setcounter{equation}{0}
\setcounter{figure}{0}
\setcounter{table}{0}
\renewcommand{\theequation}{S\arabic{equation}}
\renewcommand{\thefigure}{S\arabic{figure}}

\section{S1. Full Proof of Theorem 1}

\subsection{S1.1 Precise Statement}

\textbf{Theorem 1.} Let $H = \sum_{n=1}^N \epsilon_n |n\rangle\langle n|$ with $\epsilon_1 < \epsilon_2 < \cdots < \epsilon_N$ be a non-degenerate Hamiltonian. Let $\rho(t)$ be a differentiable trajectory of density matrices. At any time $t$ where the eigenvalues $\{r_k(t)\}$ of $\rho(t)$ are pairwise distinct, the ergotropy $\mathcal{E}(t) = \mathrm{Tr}[\rho H] - \sum_k r_k \epsilon_{\sigma(k)}$ is differentiable and satisfies
\begin{equation}
\dot{\mathcal{E}}(t) = \mathrm{Tr}[\Delta H_\rho \cdot \dot\rho],
\end{equation}
where $\Delta H_\rho = H - H_{\mathrm{pass}}^{(\rho)}$ and $H_{\mathrm{pass}}^{(\rho)} = \sum_{k=1}^N \epsilon_{\sigma(k)} |r_k\rangle\langle r_k|$.

\subsection{S1.2 Proof}

\textbf{Step 1. Differentiability of eigenvalues.}
For a Hermitian matrix family $\rho(t)$ with non-degenerate spectrum at $t_0$, the Kato--Rellich theorem guarantees that eigenvalues $r_k(t)$ and eigenprojections $|r_k(t)\rangle\langle r_k(t)|$ are analytic functions of $t$ in a neighborhood of $t_0$~\cite{Kato1966}.

\textbf{Step 2. Local constancy of the permutation.}
The permutation $\sigma$ is determined by the ordering of $\{r_k\}$. Since $r_i(t_0) \neq r_j(t_0)$ for all $i \neq j$, by continuity there exists $\delta > 0$ such that the ordering is preserved for $|t - t_0| < \delta$. Hence $\sigma$ is constant on this interval.

\textbf{Step 3. Hellmann--Feynman differentiation.}
From $\rho(t)|r_k(t)\rangle = r_k(t)|r_k(t)\rangle$, differentiate and project onto $\langle r_k|$:
\begin{equation}
\dot r_k = \langle r_k|\dot\rho|r_k\rangle.
\end{equation}

\textbf{Step 4. Assembly.}
\begin{align}
\dot{\mathcal{E}} &= \mathrm{Tr}[\dot\rho\, H] - \sum_k \dot r_k\, \epsilon_{\sigma(k)} \nonumber\\
&= \mathrm{Tr}[\dot\rho\, H] - \sum_k \langle r_k|\dot\rho|r_k\rangle\, \epsilon_{\sigma(k)} \nonumber\\
&= \mathrm{Tr}[\dot\rho\, H] - \mathrm{Tr}\!\left[\dot\rho \sum_k \epsilon_{\sigma(k)}|r_k\rangle\langle r_k|\right] \nonumber\\
&= \mathrm{Tr}\!\bigl[\dot\rho\,(H - H_{\mathrm{pass}}^{(\rho)})\bigr] = \mathrm{Tr}[\Delta H_\rho \cdot \dot\rho]. \quad \square
\end{align}

\subsection{S1.3 Regularity at Eigenvalue Crossings}

At a crossing point $t_c$ where $r_i(t_c) = r_j(t_c)$:

(a) The passive-state energy $E_{\mathrm{pass}} = \sum_k r_k \epsilon_{\sigma(k)}$ is continuous at crossings: the jump in $\sigma$ involves swapping two equal eigenvalues, contributing zero change.

(b) Therefore $\mathcal{E}(t) = \mathrm{Tr}[\rho H] - E_{\mathrm{pass}}$ is continuous.

(c) Lipschitz continuity: $|\mathcal{E}(t_1) - \mathcal{E}(t_2)| \leq 2(\epsilon_N - \epsilon_1)\|\dot\rho\|_1|t_1-t_2|$.

\textit{Proof.} Write $|\mathcal{E}(t_1)-\mathcal{E}(t_2)| \leq |\mathrm{Tr}[\rho(t_1)H] - \mathrm{Tr}[\rho(t_2)H]| + |E_{\mathrm{pass}}(t_1)-E_{\mathrm{pass}}(t_2)|$. The first term is bounded by $\|H\|_\infty \|\rho(t_1)-\rho(t_2)\|_1 \leq (\epsilon_N-\epsilon_1)\|\dot\rho\|_1|t_1-t_2|$. For the second term, $E_{\mathrm{pass}} = \sum_k r_k \epsilon_{\sigma(k)}$ with $|\epsilon_{\sigma(k)}| \leq \epsilon_N$ and $\sum_k|r_k(t_1)-r_k(t_2)| \leq \|\rho(t_1)-\rho(t_2)\|_1$ (Weyl's inequality), giving the same bound. $\square$

(d) At crossings, $\dot{\mathcal{E}}$ has finite left and right limits.

\subsection{S1.4 Qubit Verification}

For $H = \frac{\omega}{2}\sigma_z$ and $\rho = \frac{1}{2}(I + \vec{r}\cdot\vec\sigma)$, direct calculation gives $\mathcal{E} = \frac{\omega}{2}(|\vec{r}| + r_z)$ (for inverted populations). Then:
\begin{equation}
\dot{\mathcal{E}}_{\mathrm{direct}} = \frac{\omega}{2}\left(\frac{\vec{r}\cdot\dot{\vec{r}}}{|\vec{r}|} + \dot r_z\right).
\end{equation}
From Eq.~(S1): $\Delta H_\rho = \frac{\omega}{2}(\sigma_z + \hat{r}\cdot\vec\sigma)$, giving
\begin{equation}
\mathrm{Tr}[\Delta H_\rho \cdot \dot\rho] = \frac{\omega}{2}\dot r_z + \frac{\omega}{2}\frac{\vec{r}\cdot\dot{\vec{r}}}{|\vec{r}|} = \dot{\mathcal{E}}_{\mathrm{direct}}. \quad \checkmark
\end{equation}

\section{S2. Proof of Theorem 2}

\subsection{S2.1 Frequency-Dependent Effective Temperature}

Under the secular approximation, the steady state is diagonal: $\rho_{ss} = \sum_n p_n^{ss}|n\rangle\langle n|$. For each Bohr frequency $\omega_j$ connecting levels $m$ and $n$:
\begin{equation}
\frac{p_n^{ss}}{p_m^{ss}} = e^{-\beta_{\mathrm{eff}}(\omega_j)\hbar\omega_j},
\end{equation}
with $\beta_{\mathrm{eff}}(\omega_j)$ given by Eq.~(3) of the main text.

\subsection{S2.2 Strict Monotonicity}

For two baths ($M=2$), let $\alpha \equiv \alpha_1(\omega)$. Then
\begin{equation}
\frac{\partial\beta_{\mathrm{eff}}}{\partial\alpha} = \frac{e^{\beta_1\hbar\omega} - e^{\beta_2\hbar\omega}}{\hbar\omega\bigl(\alpha e^{\beta_1\hbar\omega} + (1-\alpha)e^{\beta_2\hbar\omega}\bigr)}.
\end{equation}
When $T_1 \neq T_2$, the numerator is nonzero, so $\beta_{\mathrm{eff}}$ is strictly monotone in $\alpha$. Condition (C2) then implies $\beta_{\mathrm{eff}}(\omega_a) \neq \beta_{\mathrm{eff}}(\omega_b)$, contradicting the Gibbs requirement of a single $\beta$. This proves (i). $\square$

\subsection{S2.3 Quantitative Lower Bound}

By the Cram\'{e}r--Rao inequality for exponential families:
\begin{equation}
\mathcal{N} \geq \frac{(\delta\beta)^2}{2\,I_F(\beta^*)},
\end{equation}
where $\delta\beta = \beta_{\mathrm{eff}}(\omega_a) - \beta_{\mathrm{eff}}(\omega_b)$ and $I_F(\beta) \leq (\epsilon_N - \epsilon_1)^2/4$. This gives Eq.~(4) of the main text. $\square$

\textit{Tightness of the bound.} The Cram\'{e}r--Rao bound uses only a single parameter ($\beta$) and the worst-case Fisher information. For the SSD qutrit with representative parameters ($\hbar\omega_{13}=100$~meV, $\hbar\omega_{12}=60$~meV, $T_H=300$~K, $T_C=77$~K), numerical evaluation gives $\mathcal{N}_{\mathrm{exact}} = 5.1 \times 10^{-2}$ versus $\mathcal{N}_{\mathrm{bound}} = 2.5 \times 10^{-6}$, a ratio of $\sim 2\times 10^4$. The bound remains loose by 3--4 orders of magnitude across the SSD parameter space ($\sim 10^3$ near threshold, $\sim 10^5$ far from threshold). This is typical of Cram\'{e}r--Rao bounds applied to low-dimensional discrete distributions. The bound's primary role is \emph{qualitative}: it guarantees $\mathcal{N} > 0$ under conditions (C1)--(C2) with an analytically transparent dependence on system parameters, rather than providing a tight quantitative estimate. For quantitative non-Gibbs departure, direct numerical evaluation of $\mathcal{N}$ (as performed in Sec.~S5) is recommended.

\subsection{S2.4 SSD Sufficient Condition for Population Inversion}

In the SSD topology (H$\leftrightarrow|1\rangle$--$|3\rangle$, C$\leftrightarrow|1\rangle$--$|2\rangle$):
\begin{equation}
p_3^{ss}/p_1^{ss} = e^{-\hbar\omega_{13}/k_BT_H}, \quad p_2^{ss}/p_1^{ss} = e^{-\hbar\omega_{12}/k_BT_C}.
\end{equation}
Inversion $p_3 > p_2$ requires $T_H/T_C > \omega_{13}/\omega_{12}$. $\square$

\section{S3. Thermodynamic Consistency}

\subsection{S3.1 Energy Balance}

In the SSD steady state, each transition satisfies detailed balance with its bath independently. Heat flow occurs during the replenishment phase after extraction. Per cycle:
\begin{equation}
Q_H = \mathcal{E}_{ss}/\eta, \quad Q_C = \mathcal{E}_{ss}(1/\eta - 1),
\end{equation}
satisfying $Q_H = W + Q_C$. $\checkmark$

\subsection{S3.2 Entropy Production}

The entropy production rate:
\begin{equation}
\dot\sigma = \frac{P}{\omega_{23}}\left(\frac{\omega_{12}}{T_C} - \frac{\omega_{13}}{T_H}\right).
\end{equation}
The SSD condition $T_H/T_C > \omega_{13}/\omega_{12}$ implies $\omega_{12}/T_C > \omega_{13}/T_H$, hence $\dot\sigma > 0$. $\checkmark$

\subsection{S3.3 Carnot Bound}

From $\dot\sigma \geq 0$ and $P > 0$:
\begin{equation}
\eta = \frac{\omega_{23}}{\omega_{13}} = 1 - \frac{\omega_{12}}{\omega_{13}} \leq 1 - \frac{T_C}{T_H} = \eta_{\mathrm{Carnot}}.
\end{equation}
Equality holds at the inversion threshold where $\mathcal{E}_{ss} \to 0$ (reversible, zero power). $\square$

\section{S4. Strong-Coupling Analysis}

\subsection{S4.1 Reaction-Coordinate Mapping}

Following Strasberg \textit{et al.}~\cite{Strasberg2016}, each Drude-Lorentz bath is mapped to a reaction coordinate (RC):
\begin{equation}
H_{\mathrm{ext}} = H_S + \sum_{k\in\{H,C\}} \bigl[\Omega_k b_k^\dagger b_k + g_k V_k(b_k + b_k^\dagger)\bigr],
\end{equation}
where $g_k = \sqrt{\lambda_k\Omega_k/2}$, $V_H = |1\rangle\langle3| + \mathrm{h.c.}$, $V_C = |1\rangle\langle2| + \mathrm{h.c.}$ The Hilbert space is truncated to $3 \times N_H \times N_C$ ($N_H = N_C = 6$--8). The residual environment is treated in Born-Markov approximation.

\subsection{S4.2 Numerical Protocol}

For each parameter set $(\alpha_H, \alpha_C, T_H, T_C, \omega_{13}, \omega_{23})$:
\begin{enumerate}
\item Construct $H_{\mathrm{ext}}$ and Lindblad dissipators for residual baths.
\item Solve for steady state $\tilde\rho_{ss}$ via sparse direct methods.
\item Compute $\rho_S = \mathrm{Tr}_{\mathrm{RC}}[\tilde\rho_{ss}]$.
\item Compute $\mathcal{E}(\rho_S, H_S)$ and $\mathcal{F} = \mathcal{E}^{(\mathrm{RC})}/\mathcal{E}^{(\mathrm{BM})}$.
\end{enumerate}
Convergence verified by increasing $N_{H,C}$ from 4 to 8 (all results converged to $< 1\%$).

\subsection{S4.3 Results}

Across $\alpha \in [0.01, 0.5]$:
\begin{equation}
\mathcal{F}(\alpha) \in [0.13,\, 0.21] \quad (\text{all values} < 1).
\end{equation}
Strong coupling suppresses ergotropy. The mechanism: RC hybridization delocalizes population over system+RC states, smearing the inversion. The secular condition (C1) is violated at strong coupling.

\subsection{S4.4 Analytical Estimates: Two Levels of Approximation}

We present two analytical estimates that capture different physics and give opposite predictions. Their comparison with numerics identifies which effects dominate.

\subsubsection{S4.4a. Leading-order level-shift estimate ($\mathcal{F} > 1$)}

The simplest effect of strong coupling is energy-level renormalization. At resonance ($\Omega_H = \omega_{13}$), the dressed states $|\tilde\pm\rangle = (|3,0\rangle \pm |1,1\rangle)/\sqrt{2}$ have energies $E_{\tilde\pm} = \epsilon_3 \pm g_H$. The lower dressed state $|\tilde-\rangle$ effectively shifts the energy of $|3\rangle$ downward:
\begin{equation}\label{eq:S16a}
\tilde\epsilon_3 \approx \epsilon_3 - g_H^2/\Omega_H, \quad \tilde\omega_{13} = \omega_{13} - g_H^2/(\hbar\Omega_H).
\end{equation}
This lowers the SSD inversion threshold from $T_H/T_C > \omega_{13}/\omega_{12}$ to $T_H/T_C > \tilde\omega_{13}/\omega_{12}$. The resulting enhancement factor, computed by substituting $\tilde\omega_{13}$ into Eq.~(7) of the main text while keeping $\omega_{12}$ fixed, is
\begin{equation}\label{eq:S16a-F}
\mathcal{F}^{(\mathrm{level\text{-}shift})} = \frac{\mathcal{E}_{ss}(\tilde\omega_{13})}{\mathcal{E}_{ss}(\omega_{13})} > 1 \quad (\text{for all } g_H > 0).
\end{equation}
This estimate predicts that strong coupling \emph{enhances} ergotropy by making inversion easier. It is plotted as the green curve in Fig.~2(b) of the main text.

\subsubsection{S4.4b. Complete RC trace-out ($\mathcal{F} < 1$)}

The level-shift estimate neglects a crucial effect: the RC mode enlarges the Hilbert space, diluting population across system+RC states. We now account for this.

In the dressed-state picture ($n_{\max}=1$, resonance $\Omega_H = \omega_{13}$):
\begin{equation}
|\tilde\pm\rangle = \frac{1}{\sqrt{2}}(|3,0\rangle \pm |1,1\rangle), \quad E_{\tilde\pm} = \epsilon_3 \pm g_H.
\end{equation}
The global thermal state at inverse temperature $\beta_H$ populates these dressed states as $\tilde p_\pm \propto e^{-\beta_H E_{\tilde\pm}}$. Tracing out the RC degree of freedom, the reduced population of level $|3\rangle$ is
\begin{align}
p_3^{(\mathrm{RC})} &= |\langle 3,0|\tilde+\rangle|^2 \tilde p_+ + |\langle 3,0|\tilde-\rangle|^2 \tilde p_- \nonumber\\
&= \tfrac{1}{2}(\tilde p_+ + \tilde p_-) = \tfrac{1}{2}\frac{e^{-\beta_H(\epsilon_3 - g_H)} + e^{-\beta_H(\epsilon_3 + g_H)}}{Z} \nonumber\\
&= \frac{e^{-\beta_H\epsilon_3}\cosh(\beta_H g_H)}{Z}.
\end{align}
In the weak-coupling limit ($g_H \to 0$), this reduces to the Born-Markov result $p_3^{(\mathrm{BM})} = e^{-\beta_H\epsilon_3}/Z^{(\mathrm{BM})}$. The ratio gives:
\begin{equation}
\frac{p_3^{(\mathrm{RC})}}{p_3^{(\mathrm{BM})}} = \cosh(\beta_H g_H) \cdot \frac{Z^{(\mathrm{BM})}}{Z^{(\mathrm{RC})}}.
\end{equation}
Since the RC adds an extra thermal mode, $Z^{(\mathrm{RC})} > Z^{(\mathrm{BM})}$, and the net effect on the \emph{inversion} $p_3 - p_2$ (where $p_2$ is unaffected by RC$_H$) is a reduction. Expanding to leading order in $g_H$:
\begin{equation}\label{eq:S16b}
\mathcal{F}^{(\mathrm{full\; RC})} \approx 1 - \frac{g_H^2}{\Omega_H^2}\coth\!\left(\frac{\beta_H\Omega_H}{2}\right) < 1,
\end{equation}
which is strictly less than unity for all $g_H > 0$ and all temperatures.

\subsubsection{S4.4c. Resolution: why the two estimates disagree}

The level-shift estimate (S4.4a) captures only the energy renormalization of $|3\rangle$ and predicts $\mathcal{F} > 1$. The complete trace-out (S4.4b) additionally accounts for:
\begin{itemize}
\item Population dilution: the $|3\rangle$ population is shared with the $|1,1_{\mathrm{RC}}\rangle$ component of the dressed state.
\item Partition function increase: $Z^{(\mathrm{RC})} > Z^{(\mathrm{BM})}$ due to the extra RC thermal mode.
\item Additional dissipation channels from the residual bath coupled to the RC.
\end{itemize}
These higher-order effects overwhelm the leading-order energy shift, reversing the sign of the correction. The numerical RC results (Fig.~2 of the main text, $\mathcal{F} \approx 0.1$--$0.2$) confirm that the complete estimate (S4.4b) is quantitatively correct: strong coupling suppresses, rather than enhances, ergotropy in the SSD topology. $\square$

\section{S5. Numerical Methods}

\subsection{S5.1 Implementation}

Born-Markov calculations use QuTiP (v5.0+). Jump operators:
\begin{align}
L_H^{(+)} &= \sqrt{\gamma_H(n_H+1)}\,|1\rangle\langle3|, \quad L_H^{(-)} = \sqrt{\gamma_H n_H}\,|3\rangle\langle1|, \\
L_C^{(+)} &= \sqrt{\gamma_C(n_C+1)}\,|1\rangle\langle2|, \quad L_C^{(-)} = \sqrt{\gamma_C n_C}\,|2\rangle\langle1|.
\end{align}
Steady states obtained via direct LU factorization.

\subsection{S5.2 Ergotropy Computation}

\begin{equation}
\mathcal{E}(\rho, H) = \mathrm{Tr}[\rho H] - \sum_k r_k^\downarrow \epsilon_k^\uparrow,
\end{equation}
where $\{r_k^\downarrow\}$ are eigenvalues of $\rho$ in decreasing order and $\{\epsilon_k^\uparrow\}$ are energy eigenvalues in increasing order.

\subsection{S5.3 Parameter Scan}

The scan covers:
$T_H/T_C \in \{1.5, 2, 3, 4, 5, 7, 10, 15, 20\}$ (9 values),
$\omega_{23}/\omega_{13} \in \{0.1, 0.2, \ldots, 0.9\}$ (9 values),
$\gamma_H/\gamma_C \in \{0.1, 0.3, 1, 3, 10\}$ (5 values).
Total: $9 \times 9 \times 5 = 405$ parameter sets. Fixed: $\omega_{13} = 1.0$, $\gamma_C = 0.01$.

\subsection{S5.4 Verification}

For all 405 points, analytical populations (Eq.~S8) match numerical steady states to $< 2\%$. The SSD boundary correctly predicts the sign of $\mathcal{E}$ at every point.

\begin{thebibliography}{9}
\bibitem{Kato1966} T.~Kato, \textit{Perturbation Theory for Linear Operators} (Springer, 1966), Chap.~II.
\bibitem{Pusz1978} W.~Pusz and S.~L.~Woronowicz, Commun. Math. Phys. \textbf{58}, 273 (1978).
\bibitem{Strasberg2016} P.~Strasberg, G.~Schaller, N.~Lambert, and T.~Brandes, New J. Phys. \textbf{18}, 073007 (2016).
\end{thebibliography}

\end{document}
